\section{Limitations and Future Work}
\label{sec:limitations-and-future-work}

OpenMLE provides an open path from executable environments to post-trained models and long-horizon search, but it does not yet realize the full vision of recursive self-improvement (RSI). We highlight five capability boundaries that we consider most consequential for closing this gap.

\textbf{Richer objectives for improving the improver.}
OpenMLE currently learns primarily from the measured outcome of an executed
solution. This signal reveals whether a program works and how well it scores,
but it does not fully capture whether a research direction is promising,
generalizable, robust, or worth additional computation. As a result, the
system is better equipped to optimize solutions than to judge which ideas
deserve to be pursued. A more capable improver will require objectives that
represent not only final performance, but also the quality of hypotheses,
reasoning processes, critiques, and transferable research strategies.

\textbf{Integrating evolutionary search with general coding agents.}
Our present system composes trained operators through an external evolutionary
harness. This separation makes training and search tractable, but it also
bounds the range of actions and interactions the model can initiate on its
own. Moving beyond this boundary requires combining evolutionary search with
general coding agents, bringing population-based exploration and flexible
agentic problem solving into a unified framework.

\textbf{Broader participation in AI development.}
The current environments ask an agent to improve external machine learning
artifacts. They provide a practical and verifiable testbed for meta-evolution,
but the agent participates in only a limited part of the broader AI development
process. Moving closer to recursive self-improvement requires agents to take
part in a larger share of this process, especially the improvement of language
models themselves.

\textbf{Evolving the evolutionary system.}
In OpenMLE, evolution operates primarily over candidate solutions, while the
evolutionary system itself remains largely fixed. A further step toward
recursive self-improvement is therefore to make the evolutionary system itself
an object of evolution.

\textbf{Richer use of experience in node expansion.}
Our experience-guided node-expansion framework remains a preliminary
prototype. Although each experience card preserves a broad set of
deterministic metadata about a node, the current parent-selection policy uses only three
factors that we consider especially important: solution quality,
parent-relative improvement, and method-family novelty. This design
demonstrates that structured experience can guide the allocation of search
budget, but it leaves much of the recorded evidence underutilized. Future
work could incorporate additional signals, while learning task-dependent rather than fixed
factor weights. More broadly, enabling the search policy itself to discover
which experience signals are predictive would provide a promising path from
hand-designed experience guidance toward an evolutionary system that
improves its own search behavior.
